\chapter{Introduction}
\label{chap:introduction}

\section{Motivation}
\label{sec:motivation}

Clinical practice guidelines are documents that provide evidence-based
recommendations for healthcare professionals. They may describe how to
diagnose a condition, which treatments to use, what targets to aim for, and
how patients should be monitored. These guidelines are published by professional and national organisations , which often use different document structures, terminology, and measurement units.

This variation makes it difficult to compare guidelines automatically. For
example, one organisation may report a fasting plasma glucose threshold in
mg/dL, while another reports the same type of measurement in mmol/L. A
simple text search can find the numbers, but it cannot determine whether
they describe the same clinical concept, convert the units, or preserve
which organisation reported each value.

Building a system that compares guidelines automatically is harder than it
might first appear, because existing tools for extracting structured
information from documents are usually built around one source at a time. A
parser or extraction rule written for one organisation's documents often
does not work on another organisation's documents, because the two use
different layouts, different wording, and different ways of presenting the
same kind of information. A system that is only ever tested on documents
from a single publisher may work well on those documents while failing on
anything else.

This thesis addresses the problem by representing guideline information in
a knowledge graph. A knowledge graph is a structured network in which
important items, such as organisations, treatments, conditions, and
numerical targets, are represented as nodes. Relationships between these
items are represented as edges. This structure allows information from
different guidelines to be compared and queried systematically.

As a concrete example, consider a clinician or medical educator who wants
to know whether the American Diabetes Association (ADA)\footnote{https://www.diabetes.org/}, the National Institute for Health and
Care Excellence~(NICE)\footnote{https://www.nice.org.uk/}, the International
Diabetes Federation (IDF)\footnote{https://www.idf.org/}, and Diabetes Canada\footnote{https://www.diabetes.ca/} use the same fasting plasma
glucose threshold to diagnose diabetes. Working from the source documents,
this means opening four guidelines, finding the relevant table or
paragraph in each, and converting between mg/dL and mmol/L by hand. In the knowledge graph built in this thesis, the thresholds are already converted to the same unit and linked when they represent the same normalised target, so the question can be answered with a single query, while each value still points back to the organisation and page that stated it.

The proposed system has four stages. First, it converts PDF guidelines into
machine-readable records. Second, it extracts relevant information from
those records, using a combination of pattern-based rules and a language
model, depending on how predictable the information is. Third, it
normalises terms and measurement units so that equivalent information can
be compared across organisations. Finally, it uses the resulting graph to generate and validate questions whose factual content is tied to specific, traceable statements from the source guidelines.

The instantiation evaluated in this thesis uses eight documents from four
organisations covering Type~1 and Type~2 diabetes. The four-stage
architecture is not tied to diabetes, but the specific rules used to read
the documents, identify entities, and compare values were built for this
corpus; applying the architecture to another domain would require adapting
these rules.

\section{Objectives and Contributions}
\label{sec:contributions}

The main contribution of this thesis is an end-to-end pipeline that converts diabetes clinical guidelines with different formats and terminology into a normalised knowledge graph that can be searched and compared. The graph is also used as a controlled factual source for generating questions. The pipeline converts documents into structured records, extracts relevant information, standardises entities and measurement units, compares information across guidelines, and generates and validates questions from the resulting graph. Its key
outcome is that relations between guidelines, such as two organisations
using the same diagnostic threshold once units are normalised, become
explicit and queryable in the graph, while each statement keeps its source
organisation and page. Its key outcome is that relationships between information stated in different guidelines, such as two organisations using the same diagnostic threshold once units are normalised, become explicit and queryable in the graph, while each statement keeps its source organisation and page.

\textbf{A four-stage pipeline for multi-guideline knowledge graph
construction.} The pipeline consists of an ingestion stage that converts
heterogeneous PDF documents into typed element records, an extraction stage
with rule-based and contextual components, a normalisation stage that
deduplicates entities and converts units, and a question-generation stage
that produces a validated question catalogue from the graph. The stages
have defined inputs and outputs, allowing the processing steps to be
applied and evaluated separately.

\textbf{A dual-layer extraction architecture.} The extraction stage
separates entity types between a rule-based layer and a contextual
extraction layer. Rule-based extraction handles types with predictable
lexical or structural patterns, including numerical targets and explicitly
delimited recommendations, while contextual extraction is intended for
drugs, clinical conditions, populations, and interventions that require
more contextual interpretation. Each extracted record retains its
extraction method for later analysis. The final corpus uses the implemented extraction paths described in Chapter~\ref{chap:evaluation}. The contextual extraction path was executed for the four entity categories of drugs, clinical conditions, populations, and interventions.

\textbf{Cross-guideline equivalence edges created by normalisation.} After
unit conversion, two numerical targets (for example, two blood glucose
thresholds) from different organisations that state the same measurement,
comparison, and value are linked by a directed ``equivalent'' connection.
These edges are not directly stated in any single source document. They are
created only after independently stated targets have been normalised into
a common representation and compared across organisations. In the diabetes
corpus, the final graph contains eight such connections, linking diagnostic
thresholds across ADA, IDF, and Diabetes Canada.

\textbf{A graph-anchored question generation and validation pipeline.} The
question generation stage selects graph edges as anchors and provides the
selected graph information as factual context to the language model. The
validator then checks the generated question against the graph anchor,
including answer consistency, numerical values, and guideline attribution.
This provides an automated check that the question is consistent with its
graph anchor, although it does not establish independent clinical
correctness. Questions derived from cross-organisation equivalence
connections draw on information from two guidelines simultaneously, which
no single-guideline source in this corpus can support on its
own~\cite{guo2024,bi2020}.

\section{Thesis Structure}
\label{sec:structure}

Chapter~\ref{chap:background} reviews the background concepts on which this
thesis builds: knowledge graphs and graph data models, information
extraction including named entity recognition and relation extraction,
language model grounding and hallucination, question generation, and
clinical practice guidelines.

Chapter~\ref{chap:construction} defines the four-stage pipeline formally
and describes the design decisions behind the schema, the dual-layer
extraction architecture, the normalisation functions, and the question
generation and validation approach.

Chapter~\ref{chap:evaluation} instantiates the pipeline on the diabetes
guideline corpus and reports the results: the ingestion statistics, the
final knowledge graph, the equivalence edges found by normalisation, the
question catalogue, and the outcome of a human clinical review of
30~questions by two medical doctors.

Chapter~\ref{chap:conclusion} concludes with a summary of contributions and
directions for future work.